% GENERATED FILE. Edit canonical lessons, then run npm run generate:books.
% canonical-source-hash: fnv1a64:ed0c21c3ecb3394d
% canonical-lessons: RU-C272-zaranee, RU-C272-nakanune, RU-C272-prezhde, RU-C272-vskore, RU-C272-spustya

\chapter{Beforehand and soon after}
\label{ch:ru-beforehand-and-soon-after}
\hlchaptermodality{272}

\begin{chapteropening}
\textbf{By the end of this chapter:} \emph{I can say when something happened: in advance, on the eve, before, soon after and some time later.}

Complete the last lesson of chapter 272: I can say when something happened: in advance, on the eve, before, soon after and some time later.
\end{chapteropening}

\section[\texorpdfstring{\ru{заранее} (zaráneye)}{заранее (zaráneye)}]{\ru{заранее} (zaráneye) — in advance, beforehand}

\label{lesson:RU-C272-zaranee}



\begin{quote}
Before the new one: say the Russian for immediately, then the Russian for an instant.
\end{quote}

\subsection*{You'll want to know: \ru{заранее}}
\textbf{\ru{заранее}} (\emph{zaráneye}) — \textquotedblleft{}in advance, beforehand\textquotedblright{}.

\begin{quote}
\ru{Мы} \ru{купили} \ru{билеты} \ru{заранее}.
\end{quote}

\begin{quote}
\emph{My kupíli biléty zaráneye.}
\end{quote}

\begin{quote}
We bought the tickets in advance.
\end{quote}

\emph{\ru{Заранее} \ru{спасибо}} — thanks in advance.

\begin{cousinweb}
\emph{\ru{заранее}} is built on \emph{\ru{за}-} and \emph{\ru{ранний}} \textquotedblleft{}early\textquotedblright{}, the adjective of \emph{\ru{рано}}, which is inherited from Proto-Slavic \emph{ranŭ} \textquotedblleft{}early\textquotedblright{} (EtymDB).
\end{cousinweb}

\begin{grammarlens}[title={What you've built}]
One more word for telling, explaining and describing at B1.
\end{grammarlens}

\subsection*{Your turn}
\begin{itemize}
\raggedright
  \item \emph{Say it:} \emph{zaráneye}
  \item \emph{Say it:} \emph{zaráneye}, once more
  \item \emph{Say it:} a sentence of your own with \emph{zaráneye}
  \item \emph{Recall:} say the Russian for immediately, then the Russian for an instant, then say \emph{zaráneye} again
\end{itemize}

\begin{tcolorbox}[breakable,colback=teal!4,colframe=teal!35!black,title={Before you move on}]
What does \ru{заранее} mean? (\textquotedblleft{}In advance, beforehand\textquotedblright{}.) Say it once more.
\end{tcolorbox}
\section[\texorpdfstring{\ru{накануне} (nakanúne)}{накануне (nakanúne)}]{\ru{накануне} (nakanúne) — the day before, on the eve}

\label{lesson:RU-C272-nakanune}



\begin{quote}
Before the new one: say the Russian for an instant, then the Russian for in advance.
\end{quote}

\subsection*{You'll want to know: \ru{накануне}}
\textbf{\ru{накануне}} (\emph{nakanúne}) — \textquotedblleft{}the day before, on the eve\textquotedblright{}.

\begin{quote}
\ru{Накануне} \ru{свадьбы} \ru{она} \ru{не} \ru{спала}.
\end{quote}

\begin{quote}
\emph{Nakanúne svádby oná ne spalá.}
\end{quote}

\begin{quote}
On the eve of the wedding she did not sleep.
\end{quote}

With the genitive: \emph{\ru{накануне} \ru{праздника}} — on the eve of the holiday. Alone it means \textquotedblleft{}the day before\textquotedblright{}.

\begin{cousinweb}
\emph{\ru{накануне}} is built on \emph{\ru{на}} \textquotedblleft{}on\textquotedblright{} and \emph{\ru{канун}} \textquotedblleft{}the eve of a feast\textquotedblright{}, which is usually traced to Greek \emph{kan\'{\={o}}n} \textquotedblleft{}a rule\textquotedblright{}, the name of a church service sung on the eve; EtymDB lists no parent.
\end{cousinweb}

\begin{grammarlens}[title={What you've built}]
One more word for telling, explaining and describing at B1.
\end{grammarlens}

\subsection*{Your turn}
\begin{itemize}
\raggedright
  \item \emph{Say it:} \emph{nakanúne}
  \item \emph{Say it:} \emph{nakanúne}, once more
  \item \emph{Say it:} a sentence of your own with \emph{nakanúne}
  \item \emph{Recall:} say the Russian for an instant, then the Russian for in advance, then say \emph{nakanúne} again
\end{itemize}

\begin{tcolorbox}[breakable,colback=teal!4,colframe=teal!35!black,title={Before you move on}]
What does \ru{накануне} mean? (\textquotedblleft{}The day before, on the eve\textquotedblright{}.) Say it once more.
\end{tcolorbox}
\section[\texorpdfstring{\ru{прежде} (prézhde)}{прежде (prézhde)}]{\ru{прежде} (prézhde) — before, formerly}

\label{lesson:RU-C272-prezhde}



\begin{quote}
Before the new one: say the Russian for in advance, then the Russian for the day before.
\end{quote}

\subsection*{You'll want to know: \ru{прежде}}
\textbf{\ru{прежде}} (\emph{prézhde}) — \textquotedblleft{}before, formerly\textquotedblright{}.

\begin{quote}
\ru{Прежде} \ru{здесь} \ru{был} \ru{маленький} \ru{парк}.
\end{quote}

\begin{quote}
\emph{Prézhde zdes byl málenkiy park.}
\end{quote}

\begin{quote}
Formerly there was a small park here.
\end{quote}

\emph{\ru{Прежде} \ru{чем}} + infinitive — before doing: \emph{\ru{прежде} \ru{чем} \ru{уйти}}. \emph{\ru{Прежде} \ru{всего}} — first of all.

\begin{cousinweb}
\emph{\ru{прежде}} is a Church Slavonic word, \emph{prěžde} \textquotedblleft{}before\textquotedblright{} (EtymDB), with the \emph{-\ru{ре}-} of the church language; the native Russian shape of the same root is \emph{\ru{перед}} \textquotedblleft{}in front of\textquotedblright{}.
\end{cousinweb}

\begin{grammarlens}[title={What you've built}]
One more word for telling, explaining and describing at B1.
\end{grammarlens}

\subsection*{Your turn}
\begin{itemize}
\raggedright
  \item \emph{Say it:} \emph{prézhde}
  \item \emph{Say it:} \emph{prézhde}, once more
  \item \emph{Say it:} a sentence of your own with \emph{prézhde}
  \item \emph{Recall:} say the Russian for in advance, then the Russian for the day before, then say \emph{prézhde} again
\end{itemize}

\begin{tcolorbox}[breakable,colback=teal!4,colframe=teal!35!black,title={Before you move on}]
What does \ru{прежде} mean? (\textquotedblleft{}Before, formerly\textquotedblright{}.) Say it once more.
\end{tcolorbox}
\section[\texorpdfstring{\ru{вскоре} (vskóre)}{вскоре (vskóre)}]{\ru{вскоре} (vskóre) — soon after, shortly}

\label{lesson:RU-C272-vskore}



\begin{quote}
Before the new one: say the Russian for the day before, then the Russian for before.
\end{quote}

\subsection*{You'll want to know: \ru{вскоре}}
\textbf{\ru{вскоре}} (\emph{vskóre}) — \textquotedblleft{}soon after, shortly\textquotedblright{}.

\begin{quote}
\ru{Вскоре} \ru{пошёл} \ru{снег}.
\end{quote}

\begin{quote}
\emph{Vskóre poshyól sneg.}
\end{quote}

\begin{quote}
Soon after, it began to snow.
\end{quote}

It looks back from a point in the story: soon after that. \emph{\ru{Вскоре} \ru{после} \ru{обеда}} — shortly after lunch.

\begin{cousinweb}
\emph{\ru{вскоре}} is built on \emph{\ru{в}} and \emph{\ru{скоро}} \textquotedblleft{}quickly\textquotedblright{}, from \emph{\ru{скорый}} \textquotedblleft{}fast\textquotedblright{}, which is inherited through Old East Slavic \emph{skorŭ} from Proto-Slavic \emph{skorŭ} \textquotedblleft{}fast\textquotedblright{} (EtymDB).
\end{cousinweb}

\begin{grammarlens}[title={What you've built}]
One more word for telling, explaining and describing at B1.
\end{grammarlens}

\subsection*{Your turn}
\begin{itemize}
\raggedright
  \item \emph{Say it:} \emph{vskóre}
  \item \emph{Say it:} \emph{vskóre}, once more
  \item \emph{Say it:} a sentence of your own with \emph{vskóre}
  \item \emph{Recall:} say the Russian for the day before, then the Russian for before, then say \emph{vskóre} again
\end{itemize}

\begin{tcolorbox}[breakable,colback=teal!4,colframe=teal!35!black,title={Before you move on}]
What does \ru{вскоре} mean? (\textquotedblleft{}Soon after, shortly\textquotedblright{}.) Say it once more.
\end{tcolorbox}
\section[\texorpdfstring{\ru{спустя} (spustyá)}{спустя (spustyá)}]{\ru{спустя} (spustyá) — later, after (a stretch of time)}

\label{lesson:RU-C272-spustya}



\begin{quote}
Before the new one: say the Russian for before, then the Russian for soon after.
\end{quote}

\subsection*{You'll want to know: \ru{спустя}}
\textbf{\ru{спустя}} (\emph{spustyá}) — \textquotedblleft{}later, after (a stretch of time)\textquotedblright{}.

\begin{quote}
\ru{Спустя} \ru{два} \ru{года} \ru{они} \ru{снова} \ru{встретились}.
\end{quote}

\begin{quote}
\emph{Spustyá dva góda oní snóva vstrétilis.}
\end{quote}

\begin{quote}
Two years later they met again.
\end{quote}

With the accusative of time, before or after it: \emph{\ru{спустя} \ru{неделю}} or \emph{\ru{неделю} \ru{спустя}} — a week later.

\begin{cousinweb}
\emph{\ru{спустя}} is an old adverbial participle of \emph{\ru{спустить}} \textquotedblleft{}to let down\textquotedblright{}, built on \emph{\ru{пустить}} \textquotedblleft{}to let go\textquotedblright{}, inherited from Proto-Slavic \emph{pustiti} (EtymDB): time let go by.
\end{cousinweb}

\begin{grammarlens}[title={What you've built}]
One more word for telling, explaining and describing at B1.
\end{grammarlens}

\subsection*{Your turn}
\begin{itemize}
\raggedright
  \item \emph{Say it:} \emph{spustyá}
  \item \emph{Say it:} \emph{spustyá}, once more
  \item \emph{Say it:} a sentence of your own with \emph{spustyá}
  \item \emph{Recall:} say the Russian for before, then the Russian for soon after, then say \emph{spustyá} again
\end{itemize}

\begin{tcolorbox}[breakable,colback=teal!4,colframe=teal!35!black,title={Before you move on}]
What does \ru{спустя} mean? (\textquotedblleft{}Later, after (a stretch of time)\textquotedblright{}.) Say it once more.
\end{tcolorbox}
